% GENERATED FILE. Edit canonical lessons, then run npm run generate:books.
% canonical-source-hash: fnv1a64:620588606ad21968
% canonical-lessons: GE-C66-singen, GE-C66-reisen, GE-C66-zeigen, GE-C66-bringen, GE-C66-vergessen

\chapter{To sing, To travel, To show, To bring, To forget}
\label{ch:ge-to-sing-to-travel-to-show-to-bring-to-forget}
\hlchaptermodality{66}

\begin{chapteropening}
\textbf{By the end of this chapter:} \emph{I can say to sing, to travel, to show, to bring and to forget.}

Complete the last lesson of chapter 66: I can say to sing, to travel, to show, to bring and to forget.
\end{chapteropening}

\section[singen]{singen — to sing}

\label{lesson:GE-C66-singen}



\begin{quote}
Before the new one: say the German for to laugh, then the German for to cry.
\end{quote}

\subsection*{You'll want to know: singen}
\textbf{singen} — \textquotedblleft{}to sing\textquotedblright{}.

\begin{cousinweb}
\textbf{singen} $\leftarrow$ Old High German \emph{singan} $\leftarrow$ Proto-Germanic \emph{*singwaną} — inherited, the same verb as English \emph{sing}. (EtymDB 2.1)
\end{cousinweb}

\begin{grammarlens}[title={What you've built}]
One more word for this chapter.
\end{grammarlens}

\subsection*{Your turn}
\begin{itemize}
\raggedright
  \item \emph{Say it:} \emph{singen}
  \item \emph{Say it:} \emph{singen}, once more
  \item \emph{Say it:} say \emph{singen}
  \item \emph{Recall:} say the German for to laugh, then the German for to cry, then say \emph{singen} again
\end{itemize}

\begin{tcolorbox}[breakable,colback=teal!4,colframe=teal!35!black,title={Before you move on}]
What does singen mean? (\textquotedblleft{}To sing\textquotedblright{}.) Say it once more.
\end{tcolorbox}
\section[reisen]{reisen — to travel}

\label{lesson:GE-C66-reisen}



\begin{quote}
Before the new one: say the German for to cry, then the German for to sing.
\end{quote}

\subsection*{You'll want to know: reisen}
\textbf{reisen} — \textquotedblleft{}to travel\textquotedblright{}.

\begin{cousinweb}
\textbf{reisen} $\leftarrow$ Old High German \emph{reisōn}, \textquotedblleft{}to set out on a journey\textquotedblright{} $\leftarrow$ Proto-Germanic \emph{*raisōną} — inherited. (EtymDB 2.1) It is related to English \emph{rise}: a journey began by getting up and going.
\end{cousinweb}

\begin{grammarlens}[title={What you've built}]
One more word for this chapter.
\end{grammarlens}

\subsection*{Your turn}
\begin{itemize}
\raggedright
  \item \emph{Say it:} \emph{reisen}
  \item \emph{Say it:} \emph{reisen}, once more
  \item \emph{Say it:} say \emph{reisen}
  \item \emph{Recall:} say the German for to cry, then the German for to sing, then say \emph{reisen} again
\end{itemize}

\begin{tcolorbox}[breakable,colback=teal!4,colframe=teal!35!black,title={Before you move on}]
What does reisen mean? (\textquotedblleft{}To travel\textquotedblright{}.) Say it once more.
\end{tcolorbox}
\section[zeigen]{zeigen — to show}

\label{lesson:GE-C66-zeigen}



\begin{quote}
Before the new one: say the German for to sing, then the German for to travel.
\end{quote}

\subsection*{You'll want to know: zeigen}
\textbf{zeigen} — \textquotedblleft{}to show\textquotedblright{}.

\begin{cousinweb}
\textbf{zeigen} $\leftarrow$ Middle High German \emph{zeigen} $\leftarrow$ Old High German \emph{zeigōn} — inherited. EtymDB 2.1 lists it beside English \emph{teach}: both grew from an old root meaning \textquotedblleft{}to point out\textquotedblright{}.
\end{cousinweb}

\begin{grammarlens}[title={What you've built}]
One more word for this chapter.
\end{grammarlens}

\subsection*{Your turn}
\begin{itemize}
\raggedright
  \item \emph{Say it:} \emph{zeigen}
  \item \emph{Say it:} \emph{zeigen}, once more
  \item \emph{Say it:} say \emph{zeigen}
  \item \emph{Recall:} say the German for to sing, then the German for to travel, then say \emph{zeigen} again
\end{itemize}

\begin{tcolorbox}[breakable,colback=teal!4,colframe=teal!35!black,title={Before you move on}]
What does zeigen mean? (\textquotedblleft{}To show\textquotedblright{}.) Say it once more.
\end{tcolorbox}
\section[bringen]{bringen — to bring}

\label{lesson:GE-C66-bringen}



\begin{quote}
Before the new one: say the German for to travel, then the German for to show.
\end{quote}

\subsection*{You'll want to know: bringen}
\textbf{bringen} — \textquotedblleft{}to bring\textquotedblright{}.

\begin{cousinweb}
\textbf{bringen} $\leftarrow$ Old High German \emph{bringan} $\leftarrow$ Proto-Germanic \emph{*bringaną} — inherited, the same verb as English \emph{bring}. (EtymDB 2.1)
\end{cousinweb}

\begin{grammarlens}[title={What you've built}]
One more word for this chapter.
\end{grammarlens}

\subsection*{Your turn}
\begin{itemize}
\raggedright
  \item \emph{Say it:} \emph{bringen}
  \item \emph{Say it:} \emph{bringen}, once more
  \item \emph{Say it:} say \emph{bringen}
  \item \emph{Recall:} say the German for to travel, then the German for to show, then say \emph{bringen} again
\end{itemize}

\begin{tcolorbox}[breakable,colback=teal!4,colframe=teal!35!black,title={Before you move on}]
What does bringen mean? (\textquotedblleft{}To bring\textquotedblright{}.) Say it once more.
\end{tcolorbox}
\section[vergessen]{vergessen — to forget}

\label{lesson:GE-C66-vergessen}



\begin{quote}
Before the new one: say the German for to show, then the German for to bring.
\end{quote}

\subsection*{You'll want to know: vergessen}
\textbf{vergessen} — \textquotedblleft{}to forget\textquotedblright{}.

\begin{cousinweb}
\textbf{vergessen} $\leftarrow$ Old High German \emph{firgezzan} (EtymDB 2.1) — inherited, and built like English \emph{forget}: \emph{ver-}, \textquotedblleft{}away\textquotedblright{}, + an old verb meaning \textquotedblleft{}to get\textquotedblright{}.
\end{cousinweb}

\begin{grammarlens}[title={What you've built}]
One more word for this chapter.
\end{grammarlens}

\subsection*{Your turn}
\begin{itemize}
\raggedright
  \item \emph{Say it:} \emph{vergessen}
  \item \emph{Say it:} \emph{vergessen}, once more
  \item \emph{Say it:} say \emph{vergessen}
  \item \emph{Recall:} say the German for to show, then the German for to bring, then say \emph{vergessen} again
\end{itemize}

\begin{tcolorbox}[breakable,colback=teal!4,colframe=teal!35!black,title={Before you move on}]
What does vergessen mean? (\textquotedblleft{}To forget\textquotedblright{}.) Say it once more.
\end{tcolorbox}
